\chapter{General Conclusions}\label{conclusions}
\chapterquote{We don't get to choose how we start in this life. Real 'greatness' is what you do with the hand you're dealt.}{Victor Sullivan, Uncharted 3}

\newpage
This doctoral thesis presents a comprehensive computational investigation into the biophysical principles that govern the interaction, folding, and selectivity of antimicrobial peptides (AMPs) at heterogeneous lipid interfaces. By integrating unbiased molecular dynamics, well-tempered metadynamics, and emerging quantum-compatible algorithms, this work systematically addresses the fundamental limitations of empirical trial-and-error approaches in peptide discovery. The findings provide a rationally grounded, physics-based framework that spans from atomistic mechanistic characterization to the predictive, environment-conditioned design of membrane-active therapeutics.

The first major set of results, presented in Section \ref{chap1}, focused on unraveling the mechanistic drivers of peptide-membrane recognition at the atomic level[cite: 3]. Through extensive molecular dynamics and metadynamics simulations, the interactions of five distinct AMPs (magainin-2, pleurocidin, CM15, LL37, and clavanin) were evaluated against minimalist lipid bilayers carefully chosen to represent healthy mammalian (POPC), bacterial (POPE/POPG), and cancerous (CHOL/DOPE/DOPC/DPSM/DOPS) cell membranes[cite: 3]. A critical conclusion drawn from these simulations is the decisive role of the peptide's amphipathicity, specifically governed by the transverse component of its hydrophobic dipole moment[cite: 3]. The structural analysis demonstrates that this specific descriptor physically dictates the orientation of the peptide at the membrane interface, driving the molecule to embed itself parallel to the lipid surface with a spin angle that points the hydrophobic moment directly toward the hydrophobic core of the bilayer[cite: 3]. Furthermore, this study confirmed that AMPs achieve therapeutic efficacy by displaying a measurable thermodynamic preference for anionic lipid patterns prevalent in pathological cells over neutral mammalian envelopes, with peptides like LL37 exhibiting the highest specificity margins[cite: 3].

Crucially, the computational pipeline established in Section \ref{chap1} represents a methodological milestone that goes far beyond the analysis of these isolated sequences. By combining unbiased simulations with metadynamics driven by carefully selected Collective Variables (CVs)---namely the distance to the bilayer center, the cosine of the tilt angle, and the spin angle---this framework standardizes the extraction of 1D Potentials of Mean Force (PMF) and 2D Free Energy Surfaces (FES)[cite: 3]. Consequently, this rigorous protocol opens the door to the systematic, large-scale generation of biophysical datasets[cite: 3]. Because this methodology calculates standard Gibbs energies of transfer and precise structural descriptors based on statistical mechanics rather than empirical approximations, it provides the exact architecture required to construct massive, highly reliable training sets. These future datasets will be instrumental for training Machine Learning (ML) models aimed at the automated, high-throughput prediction of AMP selectivity and the design of novel sequence candidates.

To address the severe computational bottlenecks associated with exhaustive conformational sampling and the Levinthal paradox, this research explored the integration of quantum computing workflows into structural biology, as detailed in Section \ref{chap2}[cite: 2]. Recognizing that standard molecular dynamics often struggles to simulate complete folding transitions across media of varying polarities, this section successfully implemented classical simulations on quantum devices utilizing the Variational Quantum Eigensolver (VQE) algorithm[cite: 2]. 

A major methodological advancement in Section \ref{chap2} was the adaptation of a tetrahedral lattice model and a Miyazawa-Jernigan (MJ) pairwise potential modified to explicitly account for interface-driven peptide folding[cite: 2]. By introducing a 7th-degree polynomial approximation of the sign function, the model successfully simulated a smooth, continuous transition between hydrophilic and hydrophobic phases without requiring auxiliary qubits[cite: 2]. Despite the inherent hardware limitations of the current Noisy Intermediate-Scale Quantum (NISQ) era, this quantum-accelerated workflow accurately captured the sequence-dependent folding patterns of 10-amino-acid polar, nonpolar, and amphipathic peptides using only 22 qubits, proving to be a highly resource-efficient approach for evaluating interfacial biomolecular behavior[cite: 2].

Finally, the theoretical models and physical insights developed throughout the first two sections culminated in the creation of \textit{Helix Design Studio}, presented in Section \ref{chap3}[cite: 1]. \textit{Helix Design Studio} is an open-source web application engineered for the environment-conditioned design of $\alpha$-helical peptides[cite: 1]. Unlike standard generative models trained on water-soluble globular proteins, this platform implements a parametric physical Hamiltonian that explicitly computes the energetic interplay of six physical terms: local helix propensities, short-range neighbor interactions, pairwise contacts, solvent polarity partitioning, membrane charge surface coupling, and electrostatic screening[cite: 1]. 

To resolve thermodynamic inconsistency and gradient bias across varying sequence lengths, the platform standardizes all energy contributions into dimensionless Z-scores relative to randomly generated decoy ensembles[cite: 1]. This normalization ensures that optimization is driven by physical stability rather than numerical scale[cite: 1]. The deployment of this framework enables multi-objective optimization, successfully generating candidate sequences with tunable environmental specificity through one-vs-rest and cross-environment design protocols, reliably discriminating between experimentally characterized water-soluble peptides, transmembrane models (WALP/GWALP), and canonical AMPs[cite: 1].

Summing up, the main contributions of this thesis can be gathered as follows:
\begin{itemize}
    \item The biophysical characterization of the transverse hydrophobic dipole moment as the primary structural director for interfacial alignment, spin-angle orientation, and membrane anchoring (Section \ref{chap1})[cite: 3].
    \item The establishment of a robust, highly reproducible metadynamics protocol that provides the mechanistic methodology necessary for the systematic generation of high-quality, physics-based datasets to feed future Machine Learning applications (Section \ref{chap1})[cite: 3].
    \item The successful deployment of resource-efficient quantum algorithms (VQE) to model sequence-specific peptide folding across continuous lipid-water interfaces without increasing qubit overhead (Section \ref{chap2})[cite: 2].
    \item The development, validation, and public release of \textit{Helix Design Studio}, providing the scientific community with an accessible, physics-based computational workbench for the rational, environment-specific engineering of next-generation membrane-targeted therapies (Section \ref{chap3})[cite: 1].
\end{itemize}